\documentclass[11pt]{article}
\usepackage[T1]{fontenc}
\usepackage{lmodern,amsmath,amssymb,amsthm,booktabs,microtype}
\usepackage[margin=.94in,top=.78in,bottom=.82in]{geometry}
\usepackage[hidelinks]{hyperref}
\newtheorem{theorem}{Theorem}
\newtheorem{lemma}[theorem]{Lemma}
\newcommand{\ph}{\varphi}
\newcommand{\fl}[1]{\lfloor #1\rfloor}
\title{Knuth multiplication on the rows of the Wythoff array}
\author{Blake Stone mathematics repository\\[-2pt]\small Proof developed with OpenAI Codex}
\date{}
\begin{document}
\maketitle
\begin{abstract}
We prove Peter Munn's conjecture that taking the Wythoff row of a positive integer carries
Knuth's Fibonacci product to the multiplication of extended Wythoff rows. Comparing a
quadratic irrational with two thresholds gives both the output row and the exact column.
We also determine the common divisor in a construction from Horadam's Fibonacci triples.
A primitive row of determinant $121$ gives a counterexample to a proposed squarefree-part rule.
\end{abstract}

\section{Rows, columns and products}
Let $F_0=0$, $F_1=1$, $F_{i+2}=F_{i+1}+F_i$, and $\ph=(1+\sqrt5)/2$.
If $x=\sum_{i\in I}F_i$ and $y=\sum_{j\in J}F_j$ are the Zeckendorf expansions
with indices at least $2$, Knuth's product is $x\circ y=\sum_{i\in I,j\in J}F_{i+j}$.
Set $T(a,b)=(a+b,a)$ and $B_r=(\fl{(r+1)\ph},r)$ for $r\ge0$. Both ordinary row and
column numbers start at zero. The entry $W(r,i)$ is the second coordinate of
$T^{i+2}B_r$. The Wythoff array contains each positive integer exactly once
\cite{kim}. Write $R(x)$ for its row.

Munn uses the orbit of $B_r$ under all integral powers of $T$ as the extended row.
For pairs put
\[
 (a,b)*(c,d)=(ac+bd,\ ad+bc-bd).
\]
The row product $A(r,s)$ is the index of the orbit containing $B_r*B_s$.
Munn asks whether $R(x\circ y)=A(R(x),R(y))$ in OEIS A357097 \cite{munn}.

\begin{theorem}\label{main}
Let $C=B_r*B_s=(J,K)$, let $D=J-\ph K$ and $\beta=\ph-1$.
For $x=W(r,i)$ and $y=W(s,j)$ the location of $x\circ y$ is
\[
\begin{array}{c|c|c}
\text{condition}& R(x\circ y)=A(r,s) & \text{column of }x\circ y\\\hline
D<\beta & 2K-J & i+j+4\\
\beta<D<\ph & K & i+j+2\\
D>\ph & J+K & i+j
\end{array}
\]
The three cases are equivalently $J<L$, $J=L$, and $J>L$, where
$L=\fl{(K+1)\ph}$. The endpoints cannot occur. Each case occurs.
\end{theorem}
\noindent
For example, taking both inputs in column zero, the rows $(2,2)$, $(1,2)$
and $(1,1)$ give $6\circ6=84$ in row $4$, column $4$,
$4\circ6=58$ in row $8$, column $2$, and $4\circ4=40$ in row $15$, column $0$.

\newpage
\section{A canonical pair for each integer}
The two embeddings are
\[
 U(a,b)=a+\beta b,\qquad D(a,b)=a-\ph b.
\]
Each sends $*$ to ordinary multiplication. They send $T$ to multiplication by
$\ph$ and $-\beta$, respectively, using $\beta\ph=1$ and $\beta^2=1-\beta$.
Let $P(n)=(\fl{(n+1)\ph}-1,n)$. For $n\ge0$ irrationality gives
\begin{equation}\label{strip}
    -\beta^2<D(P(n))<\beta.
\end{equation}
Conversely, the unique integer $a$ in $-\beta^2<a-\ph n<\beta$ is
$\fl{(n+1)\ph}-1$. The interval has length one. If two pairs satisfy
\eqref{strip}, the absolute value of the internal coordinate of their product is
less than $\beta^2$, since $\beta^2<\beta$. Thus the product also satisfies
\eqref{strip}.

\begin{lemma}\label{knuth}
For all positive integers $x,y$, $P(x)*P(y)=P(x\circ y)$.
\end{lemma}
\begin{proof}
Write $V_i=(F_{i+1},F_i)=T^i(1,0)$. The identities
$T(a)*b=T(a*b)=a*T(b)$ give $V_i*V_j=V_{i+j}$.
We show that $P(x)=\sum_{i\in I}V_i$. Its second coordinate is $x$.
Its first is determined by \eqref{strip}, since
$D(\sum V_i)=\sum_{i\in I}(-\beta)^i$. For every finite binary word,
including the empty word,
\[
   -1< t=\sum_{k=0}^{m-1}d_k(-\beta)^k <\ph.
\]
Indeed $t=d_0-\beta t'$ for the remaining word, and both $-\beta t'$ and
$1-\beta t'$ lie in $(-1,\ph)$. Multiplication by $\beta^2$ gives
\eqref{strip} for words starting at index $2$. This argument allows consecutive
one-digits.
By bilinearity the second coordinate of $P(x)*P(y)$ is exactly Knuth's double sum.
The product satisfies \eqref{strip}, which determines its first coordinate.
\end{proof}

The first coordinate of $P(n)$ is also
$2n-\fl{(n+1)/\ph^2}$, since $\ph+\ph^{-2}=2$.
The second coordinate of $P(x)*P(y)$ therefore recovers Lunnon's formula for
$x\circ y$ in A101330 \cite{knuth}.

\section{Finding the output row}
The anchors have $U(B_r)>0$ and $\beta<D(B_r)<\ph$. Conversely, every integral
pair $(J,K)$ with $U>0$ and $\beta<D<\ph$ is $B_K$. Indeed
$K=(U-D)/(\ph+\beta)>-1$, and
$\ph(K+1)-1<J<\ph(K+1)$ determines $J$. Each $T$-orbit has at most one
anchor. Odd shifts change the sign of $D$, and a nonzero even shift scales it
by a power of $\beta^2$ or $\ph^2$. The endpoints of the anchor interval
have ratio $\ph/\beta=\ph^2$.

\newpage
\begin{proof}[Proof of Theorem~\ref{main}]
Here $D(C)\in(\beta^2,\ph^2)$ and $U(C)>0$. If $D(C)=\beta$ then
$J+1=\ph(K+1)$, forcing $(J,K)=(-1,-1)$. If $D(C)=\ph$, then
$J=\ph(K+1)$, forcing $(J,K)=(0,-1)$. Both have $U<0$.
Thus the unique anchor of the product orbit is
\[
 B_{A(r,s)}=
 \begin{cases}
 T^{-2} C=(J-K,\,2K-J),&D(C)<\beta,\\
 C,&\beta<D(C)<\ph,\\
 T^2 C=(2J+K,\,J+K),&D(C)>\ph.
 \end{cases}
\]
In the first case the new $D$ is $\ph^2D(C)\in(1,\ph)$. In the last it is
$\beta^2D(C)\in(\beta,1)$. The physical coordinate stays positive.
This verifies Munn's orbit definition for arbitrary starting columns, since shifting
either factor shifts their product by the same amount.

For $i\ge0$, $T^{i+2}B_r=P(W(r,i))$. At $i=0$ its internal coordinate is
$\beta^2D(B_r)\in(0,\beta)$. Each further shift multiplies this by
$-\beta$, preserving \eqref{strip}.
Now multiply the pairs at the two entries. Lemma~\ref{knuth} gives
$P(x\circ y)=T^{i+j+4}C$. Substitution of the three anchors gives the indicated
columns. The examples on page 1 establish occurrence.
\end{proof}

\section{The common divisor in Horadam's construction}
Kenneth J. Ramsey's comment in A022344 proposes a squarefree-part formula for
the determinant of a sequence of hypotenuses interlaced with their differences
\cite{ramsey}. Write $x,y,x+y,x+2y$ for four consecutive terms in a row.
The hypotenuse is $H=x^2+2xy+2y^2$, with legs $2y(x+y)$ and $x(x+2y)$.
The second leg corrects the factor $xy$ in the comment. Between this hypotenuse
and the next insert
\[
     K=H(y,x+y)-H(x,y)=(x+y)(x+3y).
\]
Under $(x,y)\mapsto(y,x+y)$ the pair $(H,K)$ becomes $(H+K,H+2K)$, so the
interleaved sequence obeys the Fibonacci recurrence. Put $\Delta(x,y)=y^2-xy-x^2$.
Expansion gives
\begin{equation}\label{norm}
           \Delta(H,K)=-\Delta(x,y)^2.
\end{equation}

\begin{theorem}\label{gcd}
For relatively prime nonnegative $x,y$, $\gcd(H,K)=\gcd(5,x+3y)=\gcd(5,|\Delta(x,y)|)$.
The absolute determinant after dividing $H,K$ by their common divisor is
$\bigl(|\Delta(x,y)|/\gcd(5,|\Delta(x,y)|)\bigr)^2$.
\end{theorem}
\begin{proof}
As $H=(x+y)^2+y^2$, it is coprime to $x+y$. Any common divisor
of $H,K$ divides $x+3y$. Modulo $x+3y$ one has $H\equiv5y^2$ and $y$ is
invertible, hence the gcd is $\gcd(5,x+3y)$.
Modulo five, $\Delta(x,y)\equiv-(x+3y)^2$. Equation~\eqref{norm} finishes the proof.
\end{proof}

\newpage
For arbitrary positive $x,y$, first divide the input pair by $d=\gcd(x,y)$.
Both $H$ and $K$ scale by $d^2$. Thus Theorem~\ref{gcd}, applied to the primitive
input pair, gives the determinant after normalization for every row.

\paragraph{Row 45.}
Row $45$ has anchor $(74,45)$ and begins $119,193,312,505,\ldots$.
Here $\Delta(119,193)=121$ and the squarefree part after removing square
factors is $1$. The first two interleaved terms are $134593,217776$.
Their gcd is $1$ and the absolute determinant is $14641=121^2$. The proposed
squarefree-part formula gives $1$. Reading the squarefree portion as the product
of distinct prime divisors would instead predict $11^2=121$.
The multiplicativity of the original J determinant under Munn's row product
was already recorded in A357097 and A022344 \cite{munn,ramsey}.

\paragraph{Verification.}
The Lean files prove the digit identity for arbitrary finite binary words,
the exact row and column theorem for all ordinary array indices, the primitive gcd
identity, the determinant identity, and the numerical counterexample.
The independent verifier uses integer square roots, direct Zeckendorf expansions,
and an array generated from its row recurrences. It also checks the determinant
formula after normalization for imprimitive input pairs.

\begin{thebibliography}{9}
\bibitem{kim} C. Kimberling, The Zeckendorf array equals the Wythoff array,
\emph{Fibonacci Quarterly} 33 (1995), 3--8.
\url{https://www.fq.math.ca/Scanned/33-1/kimberling.pdf}.
\bibitem{munn} P. Munn, Multiplication table for the rows of the extended Wythoff array,
OEIS \href{https://oeis.org/A357097}{A357097} (2022).
For the row map see \href{https://oeis.org/A019586}{A019586}.
\bibitem{knuth} N. J. A. Sloane, Knuth's Fibonacci product, OEIS
\href{https://oeis.org/A101330}{A101330}, including F. Lunnon's formula (2008).
\bibitem{ramsey} A. F. Horadam, Fibonacci number triples, \emph{Amer. Math.
Monthly} 68 (1961), 751--753. K. J. Ramsey's entry comment (2007) appears
in OEIS \href{https://oeis.org/A022344}{A022344}.
\end{thebibliography}
\end{document}
